\documentclass{beamer}

\usepackage{tikz}
\usetikzlibrary{positioning,matrix, arrows.meta}
\usetikzlibrary{graphdrawing.trees}
\usetikzlibrary{external}
\tikzexternalize % activate!
\usepackage{color}
\newcommand{\vet}[1]{\foreach \num in {#1}{\el{\num}}}
\newcommand{\vetempty}[1]{\foreach \num in {#1}{\elempty{\num}}}
\newcommand{\el}[1]{\tikz{\node[font=\smaller, draw]{#1}}}
\newcommand{\elempty}[1]{\tikz{\node[font=\small, minimum size=1cm, draw=white, fill=white]{#1}}}
\newcommand{\bl}{\color{blue}}
\newcommand{\gr}{\color{gray}}
\newcommand{\re}{\color{red}}
\newcommand{\tcb}[1]{\textcolor{blue}{#1}}
\newcommand{\tcg}[1]{\textcolor{green}{#1}}
\newcommand{\tcr}[1]{\textcolor{red}{#1}}
\newcommand{\ar}{\tikz{\draw [->] (0,0) -- (0,0.4)}}
\newcommand{\nd}{\color{white!0}.}

\usepackage{beamerthemeSingapore}
\usepackage{graphicx}
%\usepackage{here}
\usepackage{url}
\usepackage[brazil]{babel}
\usepackage[T1]{fontenc}
\usepackage[utf8]{inputenc}
%\usepackage[tight,raggedright]{subfigure}
\usepackage{subfigure}
\usepackage[alf]{abntex2cite}
\usepackage{listings}
\usepackage{minted}
\usepackage{algorithm2e}
\usepackage{amsmath}
\usepackage{pdfpages}
% Please add the following required packages to your document preamble:
\usepackage{booktabs}
\usepackage{ragged2e}
\usepackage{etoolbox}
\usepackage{mathtools}
\usepackage{caption}
\usepackage{xmpmulti}

\tikzset{
  treenode/.style = {align=center, inner sep=0pt, text centered,
    font=\sffamily},
  arn_b/.style = {treenode, circle, black, font=\sffamily\bfseries, draw=black,
    fill=white, text width=1.5em},% arbre rouge noir, noeud noir
  arn_bl/.style = {treenode, circle, blue, font=\sffamily\bfseries, draw=black,
    fill=white, text width=1.5em},% arbre rouge noir, noeud noir
  arn_n/.style = {treenode, circle, white, font=\sffamily\bfseries, draw=black,
    fill=black, text width=1.5em},% arbre rouge noir, noeud noir
  arn_r/.style = {treenode, circle, red, draw=red, 
    text width=1.5em, very thick},% arbre rouge noir, noeud rouge
  arn_x/.style = {treenode, rectangle, draw=black,
    minimum width=0.5em, minimum height=0.5em}% arbre rouge noir, nil
}

\captionsetup[figure]{labelformat=empty,textfont=smaller}

\DeclarePairedDelimiter{\ceil}{\lceil}{\rceil}

\apptocmd{\frame}{}{\justifying}{} % Allow optional arguments after frame.

% TODO TODO TODO TODO TODO TODO TODO
\title{08 - Árvores binárias de busca}
\author{Prof. Alex Kutzke}
\date{Junho de 2021}

\begin{document}

% ==========================================
\begin{frame}
  \begin{center}
    % TODO TODO TODO TODO TODO TODO
    \large{Setor de Educação Profissional e Tecnológica}\\
    Tecnologia em Análise e Desenvolvimento de Sistemas\\
    \vspace{1cm}			
    \small{DS141 - Estruturas de Dados II}
  \end{center}
  \titlepage

  \begin{center}
    Slides adaptados do material produzido\\pelo Prof. Rodrigo Minetto
  \end{center}
\end{frame}
% ==========================================

% ==========================================
\begin{frame}[c]{Roteiro}
  \tableofcontents

\end{frame}
% ==========================================

\section{Motivação e conceito}

\begin{frame}[c]{Motivação}
  \begin{itemize}
    \item Algoritmo de busca binária:
    \begin{itemize}
      \item Bom desempenho para busca;
      \item Mas apenas para conjuntos ordenados;
      \item Necessita de estrutura de dados com indexação direta;
    \end{itemize}
    \item Inserção e remoção em vetores ordenados:
      \begin{itemize}
        \item Necessita rearranjar elementos;
      \end{itemize}
  \end{itemize}
\end{frame}

\begin{frame}[c]{Motivação}
  \begin{itemize}
    \item É possível uma estrutura de dados tal que
      \begin{itemize}
        \item Inserção eficiente;
        \item Remoção eficiente;
        \item Busca eficiente;
        \item Tamanho flexível;
      \end{itemize}
    \item ???
    \pause
    \item Sim!
    \pause
  \item \tcb{Árvores binárias de busca} são um exemplo;
  \end{itemize}
\end{frame}

\begin{frame}[c]
  \frametitle{Árvores binárias de busca}

  \begin{itemize}
    \item \tcb{Árvores binárias de busca} são árvores binárias com uma \textbf{propriedade fundamental}:
  \end{itemize}

  \center O valor associado à \textbf{raiz} é sempre \tcr{maior} do que o valor
          associado a qualquer nó da subárvore à \tcr{esquerda} e é sempre \tcr{menor} do que
          o valor associado a qualquer nó da subárvore à \tcr{direita}.
  
\end{frame}

\begin{frame}
  \frametitle{Exemplos}
      \centering
      \begin{tikzpicture}[-,level distance=1cm]
        \node (root) [arn_b] {50}
        child { 
          node [arn_b] {30} 
          child { 
            node [arn_b] {20} 
            child { node [arn_b] {10} }
            child [missing]
          } 
          child [missing]
          child { 
            node [arn_b] {40}
            child { node [arn_b] {35} }
            child { node [arn_b] {45} }
          } 
        }
        child [missing]
        child [missing]
        child [missing]
        child { 
          node [arn_b] {90} 
          child [missing]
          child { node [arn_b] {95} } 
        };
      \end{tikzpicture}
\end{frame}

\begin{frame}
  \frametitle{Exemplos}
      \centering
      \begin{tikzpicture}[-,level distance=1cm]
        \node (root) [arn_b] {e}
        child { 
          node [arn_b] {c}
          child { node [arn_b] {a} }
          child { node [arn_b] {d} }
        }
        child [missing]
        child { 
          node [arn_b] {m}
          child [missing]
          child {
            node [arn_b] {r}
            child [missing]
            child { node [arn_b] {z} }
          }
        };
      \end{tikzpicture}
\end{frame}

\begin{frame}[c]
  \frametitle{Percurso in-ordem}

  \begin{itemize}
    \item A propriedade fundamental garante que:
  \end{itemize}

  \center Quando percorrida em ordem \tcb{simétrica}, ou seja, in-ordem (E,R,D), os valores
          são encontrados em ordem crescente.

\end{frame}

\begin{frame}
  \frametitle{Percurso in-ordem: exemplo}

    In-ordem: \texttt{10,20,30,35,40,45,50,90,95}.
    \vspace{1cm}

      \centering
      \begin{tikzpicture}[-,level distance=1cm]
        \node (root) [arn_b] {50}
        child { 
          node [arn_b] {30} 
          child { 
            node [arn_b] {20} 
            child { node [arn_b] {10} }
            child [missing]
          } 
          child [missing]
          child { 
            node [arn_b] {40}
            child { node [arn_b] {35} }
            child { node [arn_b] {45} }
          } 
        }
        child [missing]
        child [missing]
        child [missing]
        child { 
          node [arn_b] {90} 
          child [missing]
          child { node [arn_b] {95} } 
        };
      \end{tikzpicture}
\end{frame}

\begin{frame}
  \frametitle{Percurso in-ordem: exemplo}

    In-ordem: \texttt{a,c,d,e,m,r,z}.
    \vspace{1cm}

      \centering
      \begin{tikzpicture}[-,level distance=1cm]
        \node (root) [arn_b] {e}
        child { 
          node [arn_b] {c}
          child { node [arn_b] {a} }
          child { node [arn_b] {d} }
        }
        child [missing]
        child { 
          node [arn_b] {m}
          child [missing]
          child {
            node [arn_b] {r}
            child [missing]
            child { node [arn_b] {z} }
          }
        };
      \end{tikzpicture}

\end{frame}

\begin{frame}[c]
  \frametitle{Percurso in-ordem}

  \begin{itemize}
    \item A propriedade fundamental \tcb{também} possibilita que:
  \end{itemize}

  \center As operações de \tcb{busca, inserção e remoção} podem ser realizadas de forma
          \textbf{eficiente}. \tcr{Como?}

  \pause

  \begin{itemize}
    \item Basta aplicar o mesmo algoritmo da busca binária para todas as operações acima.
  \end{itemize}

\end{frame}

\section{Operações}

\begin{frame}[fragile]
  \frametitle{Estrutura}

  \begin{itemize}
    \item Podemos utilizar a mesma estrutura de dados  apresentada na última aula;
    \item Árvores binária de busca são árvores binárias normais mas com uma propriedade fundamental.
  \end{itemize}

  \begin{columns}
    \begin{column}{0.4\textwidth}
      \begin{minted}[frame=lines,framesep=2mm,baselinestretch=1,fontsize=\normalsize,linenos]{c}
typedef struct arvore {
  int info;
  struct arvore *esq;
  struct arvore *dir;
} Arvore;
      \end{minted}
    \end{column}
    \begin{column}{0.4\textwidth}
      \centering
      \begin{tikzpicture}[-]
        \node at (0,-1) (root) [arn_b] {a}
        child { node [arn_b] {b} edge from parent node [above,left] {*esq} }
        child { node [arn_b] {c} edge from parent node [above,right] {*dir} };
      \end{tikzpicture}
    \end{column}
  \end{columns}
  
\end{frame}

\begin{frame}[fragile]
  \frametitle{Busca}

  \begin{itemize}
    \item A operação de \tcb{busca} explora diretamente a propriedade fundamental da árvore:
  \end{itemize}
  
  \begin{minted}[frame=lines,framesep=2mm,baselinestretch=1,fontsize=\normalsize,linenos]{c}
int buscar (Arvore *a, int v) {
  if (a == NULL) { return 0; } /*Nao achou*/
  else if (v < a->info) {
    return buscar (a->esq, v);
  }
  else if (v > a->info) {
    return buscar (a->dir, v);
  }
  else { return 1; } /*Achou*/
}
  \end{minted}

\end{frame}

\begin{frame}[fragile]
  \frametitle{Inserção}

  \begin{itemize}
    \item A operação de \tcb{inserção} precisa garantir que propriedade fundamental da árvore seja mantida;
    \item Ou seja, o novo elemento precisa ser inserido na \tcr{posição correta}; 
    \item A solução mais simples é \tcb{sempre inserir elementos folha};
    \item Para isso seguimos o seguinte algoritmo recursivo para inserir um novo elemento \textbf{x}:
    \begin{enumerate}
      \item Se $x$ é menor que a raiz, insere na subárvore da esquerda;
      \item Se $x$ é maior que a raiz, insere na subárvore da direita;
      \item Se a árvore é vazia, apenas cria um novo nó com valor $x$.
    \end{enumerate}

  \end{itemize}
\end{frame}

\begin{frame}[fragile]
  \frametitle{Inserção: implementação}
  
  \begin{minted}[frame=lines,framesep=2mm,baselinestretch=1,fontsize=\normalsize,linenos]{c}
Arvore* inserir (Arvore *a, int v) {
  if (a == NULL) {
    a = (Arvore*)malloc(sizeof(Arvore));
    a->info = v;
    a->esq = a->dir = NULL;
  }
  else if (v < a->info) {
    a->esq = inserir (a->esq, v);
  }
  else { a->dir = inserir (a->dir, v); }
  return a;
}
  \end{minted}

\end{frame}

\begin{frame}
  \frametitle{Exemplo inserção}

    \only<1-5>{Inserir o elemento 42.}

    \vspace{1cm}

    \noindent
      \centering
      \begin{tikzpicture}[
          -, 
          level/.style={sibling distance=60mm/####1},
        ]
        \node (root) [arn_b] {50}
        child { 
          node (n30) [arn_b] {30} 
          child { 
            node (n20) [arn_b] {20} 
            child { node (n10) [arn_b] {10} }
            child [missing]
          } 
          child { 
            node (n40) [arn_b] {40}
            child { node (n35) [arn_b] {35} }
            child { 
              node (n45) [arn_b] {45}
           }
          } 
        }
        child { 
          node (n90) [arn_b] {90} 
          child [missing]
          child { node (n95) [arn_b] {95} } 
        };
        \only<2>{\node [arn_r,above right=0pt of root] {42};}
        \only<3>{\node [arn_r,above left=0pt of n30] {42};}
        \only<4>{\node [arn_r,above right=0pt of n40] {42};}
        \only<5>{\node [arn_r,above right=0pt of n45] {42};}
      \end{tikzpicture}
\end{frame}

\begin{frame}
  \frametitle{Exemplo inserção}

    Elemento 42 inserido.

      \centering
      \begin{tikzpicture}[
          -, 
          level/.style={sibling distance=60mm/####1}
        ]
        \node (root) [arn_b] {50}
        child { 
          node [arn_b] {30} 
          child { 
            node [arn_b] {20} 
            child { node [arn_b] {10} }
            child [missing]
          } 
          child { 
            node [arn_b] {40}
            child { node [arn_b] {35} }
            child { 
              node [arn_b] {45}
              child { node [arn_r] {42} }
              child [missing]
            }
          } 
        }
        child { 
          node [arn_b] {90} 
          child [missing]
          child { node [arn_b] {95} } 
        };
      \end{tikzpicture}
\end{frame}

\section{Operação de remoção}

\begin{frame}
  \frametitle{Operação: remoção}

  A operação de \tcb{remoção} é mais complexa que a inserção.
  Novamente, a implementação dessa função vem da natureza recursiva da árvore:
  \begin{itemize}
    \item Se a árvore for \tcr{vazia}, nada tem de ser feito, pois o elemento não
      está presente na árvore;
    \item Se a árvore \tcr{não for vazia}:
      \begin{enumerate}
        \item Compara-se o valor armazenado no nó raiz ao valor que se
          deseja retirar da árvore;
        \item Se o valor a ser removido for \tcr{menor} do que a raiz, chama-se a
          função recursivamente para subárvore da \tcb{esquerda};
        \item Se o valor a ser removido for \tcr{maior} do que a raiz, chama-se a
          função recursivamente para subárvore da \tcb{direita};
        \item Se o valor a ser removido for \tcr{igual} a raiz, remove-se esse elemento;
      \end{enumerate}

    \item Entretanto, teremos 3 casos diferentes para esse operação do item $4$;
  \end{itemize}
\end{frame}

\begin{frame}
  \frametitle{Operação: remoção ($1^o$ caso)}

  Três situações diferentes para a remoção de um nó na árvore binária de busca binária:

  \begin{itemize}
    \item \textbf{$1^o$ caso}: Quando o nó a ser removido é uma \tcb{folha} da árvore;
    \begin{itemize}
      \item Ou seja, uma raiz sem filhos;
      \item Basta liberar a memória alocada para o elemento removido e retornar NULL para que
        a árvore pai possa atualizar seu estado;
    \end{itemize}
  \end{itemize}

\end{frame}

\begin{frame}
  \frametitle{Exemplo remoção ($1^o$ caso)}

    \only<1-6>{Remover o elemento \tcb{10}.}
    
    \noindent
    \centering
    \begin{tikzpicture}[
        -, 
        level/.style={sibling distance=60mm/####1},
      ]
      \node (root) [arn_b] {50}
      child { 
        node (n30) [arn_b] {30} 
        child { 
          node (n20) [arn_b] {20} 
          child { node (n10) [arn_b] {10} }
          child [missing]
        } 
        child { 
          node (n40) [arn_b] {40}
          child { 
            node (n35) [arn_b] {35} 
            child [missing]
            child {node (n37) [arn_b] {37} }
          }
          child { 
            node (n45) [arn_b] {45}
         }
        } 
      }
      child { 
        node (n90) [arn_b] {90} 
        child [missing]
        child { node (n95) [arn_b] {95} } 
      };
      \only<2>{\node [arn_bl,at = (root)] {50};}
      \only<3>{\node [arn_bl,at = (n30)] {30};}
      \only<4>{\node [arn_bl,at = (n20)] {20};}
      \only<5>{\node [arn_bl,at = (n10)] {10};}
      \only<6>{\node [arn_r,at = (n10)] {10};}
    \end{tikzpicture}

\end{frame}

\begin{frame}
  \frametitle{Exemplo remoção ($1^o$ caso)}

    Remover o elemento \tcb{10} - Resultado final:
    
    \noindent
    \centering
    \begin{tikzpicture}[
        -, 
        level/.style={sibling distance=60mm/####1},
      ]
      \node (root) [arn_b] {50}
      child { 
        node (n30) [arn_b] {30} 
        child { 
          node (n20) [arn_b] {20} 
          child [missing]
          child [missing]
        } 
        child { 
          node (n40) [arn_b] {40}
          child { 
            node (n35) [arn_b] {35} 
            child [missing]
            child {node (n37) [arn_b] {37} }
          }
          child { 
            node (n45) [arn_b] {45}
         }
        } 
      }
      child { 
        node (n90) [arn_b] {90} 
        child [missing]
        child { node (n95) [arn_b] {95} } 
      };
    \end{tikzpicture}

\end{frame}

\begin{frame}
  \frametitle{Operação: remoção ($2^o$ caso)}

  \begin{itemize}
    \item \textbf{$2^o$ caso}: Quando o nó a ser removido possui apenas \tcb{um filho}
    \begin{itemize}
      \item Basta liberar a memória alocada para o elemento removido e retornar o ponteiro do único filho
        para que a árvore pai possa atualizar seu estado;
    \end{itemize}
  \end{itemize}

\end{frame}

\begin{frame}
  \frametitle{Exemplo remoção ($2^o$ caso)}

    \only<1-4>{Remover o elemento \tcb{90}.}
    
    \noindent
    \centering
    \begin{tikzpicture}[
        -, 
        level/.style={sibling distance=60mm/####1},
      ]
      \node (root) [arn_b] {50}
      child { 
        node (n30) [arn_b] {30} 
        child { 
          node (n20) [arn_b] {20} 
          child { node (n10) [arn_b] {10} }
          child [missing]
        } 
        child { 
          node (n40) [arn_b] {40}
          child { 
            node (n35) [arn_b] {35} 
            child [missing]
            child {node (n37) [arn_b] {37} }
          }
          child { 
            node (n45) [arn_b] {45}
         }
        } 
      }
      child { 
        node (n90) [arn_b] {90} 
        child [missing]
        child { node (n95) [arn_b] {95} } 
      };
      \only<2>{\node [arn_bl,at = (root)] {50};}
      \only<3>{\node [arn_bl,at = (n90)] {90};}
      \only<4>{\node [arn_r,at = (n90)] {90};}
    \end{tikzpicture}

\end{frame}

\begin{frame}
  \frametitle{Exemplo remoção ($2^o$ caso)}

    Remover o elemento \tcb{90} - Resultado final:
    
    \noindent
    \centering
    \begin{tikzpicture}[
        -, 
        level/.style={sibling distance=60mm/####1},
      ]
      \node (root) [arn_b] {50}
      child { 
        node (n30) [arn_b] {30} 
        child { 
          node (n20) [arn_b] {20} 
          child { node (n10) [arn_b] {10} }
          child [missing]
        } 
        child { 
          node (n40) [arn_b] {40}
          child { 
            node (n35) [arn_b] {35} 
            child [missing]
            child {node (n37) [arn_b] {37} }
          }
          child { 
            node (n45) [arn_b] {45}
         }
        } 
      }
      child { 
        node (n90) [arn_bl] {95} 
      };
    \end{tikzpicture}

\end{frame}

\begin{frame}
  \frametitle{Operação: remoção ($3^o$ caso)}

  \begin{itemize}
    \item \textbf{$3^o$ caso}: Quando o nó a ser removido possui \tcb{dois filhos}
    \begin{itemize}
      \item É necessário um procedimento mais elaborado para que a propriedade fundamental da árvore
        seja mantida após a remoção;
      \item A pergunta chave é: \textbf{Qual elemento da árvore deve ocupar o lugar daquele a ser removido?}
    \end{itemize}
  \end{itemize}
  \pause

  Resposta: Aquele que \tcb{imediatamente o precede} ou que \tcb{imediatamente o sucede} na ordenação dos elementos da árvore.

\end{frame}

\begin{frame}[c]
  \frametitle{Operação: remoção ($3^o$ caso)}

  \centering
  Escolheremos aquele que \textbf{imediatamente o precede}. 

  \vspace{1cm}

  Ou seja, mais próximo daquele que se deseja remover e que seja menor do que este.
  
  \vspace{1cm}


  Em outras palavras ...
  \pause
  
  \vspace{1cm}

  Encontrar o elemento \tcb{mais à direita} da \tcr{subárvore da esquerda} do elemento que se deseja remover.

\end{frame}

\begin{frame}[c]
  \frametitle{Operação: remoção ($3^o$ caso)}

  Assim, a remoção no terceiro caso segue os seguintes passos:

  \begin{enumerate}
    \item \tcb{Encontre} o elemento a ser removido;
    \item Vá para sua \tcb{subárvore da esquerda};
    \item Desça até o \tcb{último nó mais à direita} dessa subárvore;
    \item \tcb{Atribua} o valor desse nó ao nó onde se encontra o elemento a ser removido;
    \item \tcb{Remova} esse último nó;
  \end{enumerate}

  A remoção final será como nos casos 1 ou 2, pois ou ele será um nó folha ou um nó com apenas um filho
  (um filho à esquerda);
\end{frame}

\begin{frame}
  \frametitle{Exemplo remoção ($3^o$ caso)}

    \only<1-9>{Remover o elemento \tcb{50}}\only<1-8>{.}\only<9>{ - Resultado final:}
    
    \noindent
    \centering
    \begin{tikzpicture}[
        -, 
        level/.style={sibling distance=60mm/####1},
      ]
      \only<1-6>{
        \node (root) [arn_b] {50}
        child { 
          node (n30) [arn_b] {30} 
          child { 
            node (n20) [arn_b] {20} 
            child { node (n10) [arn_b] {10} }
            child [missing]
          } 
          child { 
            node (n40) [arn_b] {40}
            child { 
              node (n35) [arn_b] {35} 
              child [missing]
              child {node (n37) [arn_b] {37} }
            }
            child [missing]
          } 
        }
        child { 
          node (n90) [arn_b] {90} 
          child [missing]
          child { node (n95) [arn_b] {95} } 
        };
      }
      \only<7>{
        \node (n40) [arn_b] {40}
        child { 
          node (n30) [arn_b] {30} 
          child { 
            node (n20) [arn_b] {20} 
            child { node (n10) [arn_b] {10} }
            child [missing]
          } 
          child { 
            node (n50) [arn_b] {50}
            child { 
              node (n35) [arn_b] {35} 
              child [missing]
              child {node (n37) [arn_b] {37} }
            }
            child [missing]
          } 
        }
        child { 
          node (n90) [arn_b] {90} 
          child [missing]
          child { node (n95) [arn_b] {95} } 
        };
      }
      \only<8-9>{
        \node (n40) [arn_b] {40}
        child { 
          node (n30) [arn_b] {30} 
          child { 
            node (n20) [arn_b] {20} 
            child { node (n10) [arn_b] {10} }
            child [missing]
          } 
          child { 
              node (n35) [arn_b] {35} 
              child [missing]
              child {node (n37) [arn_b] {37} }
          } 
        }
        child { 
          node (n90) [arn_b] {90} 
          child [missing]
          child { node (n95) [arn_b] {95} } 
        };
      }
      \only<2>{\node [arn_bl,at = (root)] {50};}
      \only<3-6>{\node [arn_r, at = (root)] {50};}
      \only<4>{\node [arn_bl, at = (n30)] {30};}
      \only<5>{\node [arn_bl, at = (n40)] {40};}
      \only<6>{\node [arn_r, at = (n40)] {40};}
      \only<7-8>{\node [arn_r, at = (n40)] {40};}
      \only<7>{\node [arn_r, at = (n50)] {50};}
      \only<8>{\node [arn_r, at = (n35)] {35};}
      \only<8>{\node [arn_r, at = (n37)] {37};}
    \end{tikzpicture}

\end{frame}

\section{Exercícios propostos}

\begin{frame}
  \frametitle{Exercícios propostos}

  \textbf{Exercício 1} Desenhe as árvores binárias de busca
  para os seguintes números inseridos nessa ordem:

  \begin{itemize}
    \item 1, 2, 3, 4, 5, 6, 7;
    \item 7, 6, 5, 4, 3, 2, 1;
    \item 4, 6, 3, 5, 1, 7, 3;
  \end{itemize}
  
  Qual das três é melhor para buscar um dado elemento.
\end{frame}

\begin{frame}
  \frametitle{Exercícios propostos}

  \textbf{Exercício 2} Escreva um programa que produza 30
  números aleatórios entre 0 e 999. Insira esses números em
  uma árvore binária de busca e imprima os valores na árvore
  (utilizando um percurso in-ordem). Os valores estão ordenados?

  \vspace{1cm}

  \textbf{Exercício 3} Escreva um programa que produza 100000
  números aleatórios entre 0 e 99999. Insira esses números em uma
  árvore binária de busca. Procure por um valor que não existe,
  por exemplo 100000. Quanto tempo a busca levou?

\end{frame}


\begin{frame}
  \frametitle{Exercícios propostos}

  \textbf{Exercício 4} Escreva um programa que produza 100000
  números em ordem (0 até 99999). Insira esses números em uma
  árvore binária de busca. Procure por um valor que não existe
  por exemplo 100000. Quanto tempo a busca levou?

  \vspace{1cm}

  \textbf{Exercício 5} Escreva a função ``min'' que encontre
  uma chave mínima em uma árvore de busca binária. Escreva uma função
  ``max'' que encontre uma chave máxima.

\end{frame}

\begin{frame}
  \frametitle{Exercícios propostos}

  \textbf{Exercício 6} Escreva uma função para remover elementos
  de uma árvore binária de busca. Teste todos os 3 casos descritos
  na aula.

\end{frame}

% ========================================== 
% Frame para referências 
%\begin{frame}[allowframebreaks] 
%  \frametitle{Referências utilizadas na apresentação}
%  \bibliography{referencias}
%\end{frame}
% ==========================================

\end{document}
